% Generated by figures/src/appendix_tables.py from committed results; do not edit by hand.
\footnotesize\linespread{1}\selectfont
\setlength{\tabcolsep}{4pt}
\begin{tabular}{L{2.7cm}L{1.45cm}L{1.45cm}L{1.45cm}L{1.45cm}L{1.45cm}L{1.6cm}}
\toprule
\textbf{Census division} & \textbf{States} & \textbf{90.1-1989} & \textbf{90.1-2004} & \textbf{90.1-2010} & \textbf{90.1-2013} & \textbf{Largest state (\%)\textsuperscript{a}} \\
\midrule
New England & 6 & 1990 (6) & 2005 (6) & 2012 (4) & 2015 (3) & MA (49) \\
Middle Atlantic & 3 & 1991 (3) & 2007 (3) & 2013 (2) & 2016 (2) & NY (46) \\
East North Central & 5 & 1995 (5) & 2008 (5) & 2013 (4) & 2016 (3) & IL (28) \\
West North Central & 7 & 1992 (4) & 2007 (3) & 2014 (2) & 2016 (1) & MO (30) \\
South Atlantic & 9 & 1995 (9) & 2005 (9) & 2012 (8) & 2015 (2) & FL (27) \\
East South Central & 4 & 2005 (4) & 2007 (4) & 2013 (4) & 2016 (1) & TN (39) \\
West South Central & 4 & 1995 (4) & 2007 (4) & 2015 (2) & 2017 (2) & TX (70) \\
Mountain & 8 & 1995 (8) & 2007 (8) & 2013 (7) & 2016 (3) & AZ (27) \\
Pacific & 5 & 1994 (4) & 2005 (4) & 2016 (4) & 2016 (3) & CA (73) \\
\bottomrule
\end{tabular}
\par\smallskip\parbox{\textwidth}{\scriptsize First year any state of the division had the edition, a later one or its IECC equivalent in force, with the number of states that had it by 2018 in parentheses, including states placed at that edition or later by a DOE status snapshot.\newline \textsuperscript{a} The state with the most commercial energy use in 2018 (EIA State Energy Data System) and its share of the division, the weight used to draw a building's state.}
